\chapter{Unitarity and global hyperbolicity (KS Section 5)}

This chapter is the blueprint for Section~5 of Kontsevich--Segal,
\emph{Wick rotation and the positivity of energy in quantum field theory}
(arXiv:2105.10161), the Lorentzian and Wick-rotation capstone, planned for
\texttt{KontsevichSegal/WickRotation/}. The Lorentzian cobordism category
$\mathcal C_d^{\mathrm{Lor}}$ lies on the boundary of the complex category
$\mathcal C_d^{\mathbb C}$, in the same sense in which the real Lorentzian metrics
lie on the Shilov boundary of $Q_{\mathbb C}(V)$. On passing to this boundary, a
unitary field theory on $\mathcal C_d^{\mathbb C}$ extends to the globally
hyperbolic Lorentzian cobordisms, where it becomes a genuine unitary theory:
time-symmetric hypersurfaces acquire Hilbert spaces, and globally hyperbolic
cobordisms between them act by unitary operators. The globally hyperbolic
cobordisms play the role that the dense open subgroup of the Shilov boundary
plays in the representation theory of a complex semigroup.

Section~4 of the paper, ``Some analogies from representation theory'', is
expository. It motivates the framework by analogy with the representation theory
of $\mathrm{PSL}_2(\mathbb R)$, of the symplectic group with its oscillator
semigroup and metaplectic representation, and of $\mathrm{Diff}^+(S^1)$, in each
of which a distinguished class of unitary representations arises as the boundary
values of holomorphic representations of a complex semigroup by contraction
operators. It states no definitions, axioms, or theorems of its own, so there is
nothing to formalize from it. Within Section~5 itself, the two-dimensional
light-line discussion, including the ``black hole'' closed light-cycle, and the
$d=1$ quantum-mechanics walkthrough are worked illustrations of the proof
technique behind Theorem~5.2; they are not separate statements to formalize.

\section{The Lorentzian cobordism category}

\begin{definition}[The Lorentzian cobordism category $\mathcal C_d^{\mathrm{Lor}}$ and $\mathcal C_d^{\mathrm{Lor},\omega}$; KS Section 5]
  \label{def:lorentzian-cobordism-category}
  \notready
  \uses{def:cobordism-category, prop:lorentzian-boundary, prop:only-lorentzian}
  The \emph{Lorentzian cobordism category} $\mathcal C_d^{\mathrm{Lor}}$ is the
  Lorentzian analogue of the complex cobordism category
  $\mathcal C_d^{\mathbb C}$.

  \emph{Objects.} An object is a \emph{Lorentzian germ}: a germ of a $d$-manifold
  along a closed co-oriented $(d-1)$-manifold $\Sigma$, carrying a Lorentzian
  metric of signature $(d-1,1)$. Exponentiating the geodesics emanating
  perpendicularly from $\Sigma$ identifies a neighbourhood with
  $\Sigma \times (-\varepsilon, \varepsilon)$, in which the metric takes the form
  $h_t - \mathrm{d}t^2$, where $t \mapsto h_t$ is a smooth family of Riemannian
  metrics on $\Sigma$.

  \emph{Morphisms.} A morphism $\Sigma_0 \to \Sigma_1$ is a \emph{Lorentzian
  cobordism} $M \colon \Sigma_0 \leadsto \Sigma_1$, a compact $d$-manifold
  carrying a Lorentzian metric whose incoming and outgoing boundary germs are the
  objects $\Sigma_0$ and $\Sigma_1$. Two cobordisms represent the same morphism
  when there is an isometry between them restricting to the identity on the
  boundary germs, exactly as for $\mathcal C_d^{\mathbb C}$.

  \emph{Composition.} Composition is concatenation of cobordisms, glued along the
  common boundary germ. As for $\mathcal C_d^{\mathbb C}$, this is not literally a
  category: there are no identity morphisms.

  \emph{The boundary relation.} $\mathcal C_d^{\mathrm{Lor}}$ lies on the boundary
  of $\mathcal C_d^{\mathbb C}$. At each point the Lorentzian metric is a real
  bilinear form of signature $(d-1,1)$, and by
  Propositions~\ref{prop:lorentzian-boundary} and~\ref{prop:only-lorentzian} the
  real Lorentzian metrics are exactly the nondegenerate real metrics on the
  Shilov boundary of $Q_{\mathbb C}(V)$. The Lorentzian germs and cobordisms are
  therefore the boundary points of the complex germs and cobordisms of
  $\mathcal C_d^{\mathbb C}$, and $\mathcal C_d^{\mathrm{Lor}}$ is the Lorentzian
  Shilov boundary of the complex cobordism category.

  \emph{The real-analytic version.} $\mathcal C_d^{\mathrm{Lor},\omega}$ is the
  subcategory in which both the manifolds and their Lorentzian metrics are
  real-analytic. This restriction is required by the operator assignment of
  Section~5: the forthcoming Principle~5.1 and Theorem~5.2 move a cobordism $M$
  inside a complexification $M_{\mathbb C}$, a holomorphic $d$-manifold
  containing $M$ as a totally real submanifold with the metric induced from a
  holomorphic symmetric form on $TM_{\mathbb C}$, and such a complexification
  exists only when $M$ and its metric are real-analytic.

  Stating this in Lean requires Lorentzian-manifold geometry not developed here:
  Lorentzian metrics of signature $(d-1,1)$ and their light-cone structure, the
  geodesic-normal form $h_t - \mathrm{d}t^2$ of a germ, isometries of Lorentzian
  cobordisms and their gluing, and, for $\mathcal C_d^{\mathrm{Lor},\omega}$, the
  real-analytic structure together with the complexification $M_{\mathbb C}$ into
  which a real-analytic cobordism embeds.
\end{definition}

\section{Global hyperbolicity}

\begin{definition}[Global hyperbolicity and the subcategory $\mathcal C_d^{\mathrm{gh}}$; KS Section 5]
  \label{def:globally-hyperbolic}
  \notready
  \uses{def:lorentzian-cobordism-category, def:allowable}
  Among the Lorentzian cobordisms of $\mathcal C_d^{\mathrm{Lor}}$ we single out the
  \emph{globally hyperbolic} ones.

  \emph{The general-relativity definition.} In general relativity a Lorentzian
  cobordism $M \colon \Sigma_0 \leadsto \Sigma_1$ between Riemannian manifolds is
  \emph{globally hyperbolic} if every maximally extended timelike geodesic in $M$
  travels from $\Sigma_0$ to $\Sigma_1$. Such an $M$ is necessarily diffeomorphic to
  $\Sigma_0 \times [0,1]$. It is exactly for globally hyperbolic $M$ that the Cauchy
  problem for the wave equation on $M$ is soluble.

  \emph{The compact definition adopted here.} Kontsevich and Segal consider only
  \emph{compact} cobordisms, which are not the usual setting of relativity theory.
  In the compact case they take global hyperbolicity to be the existence of a smooth
  \emph{time-function} $t \colon M \to [0,1]$ whose gradient, with respect to the
  Lorentzian metric, lies everywhere in the positive light-cone. Such a $t$ is then a
  fibration with Riemannian fibres, the time-slices $t^{-1}(s)$, and following the
  orthogonal trajectories to the time-slices produces a diffeomorphism
  $M \to \Sigma_0 \times [0,1]$. The two characterizations agree in the compact
  setting in which the paper works: each yields the product structure
  $M \cong \Sigma_0 \times [0,1]$, and the time-function form is the definition
  adopted.

  \emph{The globally hyperbolic subcategory.} The globally hyperbolic cobordisms are
  closed under concatenation, so the objects of $\mathcal C_d^{\mathrm{Lor}}$ together
  with the globally hyperbolic cobordisms between them form a subcategory
  $\mathcal C_d^{\mathrm{gh}} \subseteq \mathcal C_d^{\mathrm{Lor}}$. The existence of
  a time-function on a compact Lorentzian cobordism is an \emph{open} condition, so
  $\mathcal C_d^{\mathrm{gh}}$ is an \emph{open} subcategory of
  $\mathcal C_d^{\mathrm{Lor}}$. It is the analogue of the dense open subgroup of the
  Shilov boundary of a complex semigroup, the real Lie group to which the holomorphic
  contraction representations of Section~4 extend as genuine unitary operators. It
  should play that role here, although the theorem obtained below, Theorem~5.2, is
  weaker.

  \emph{The real-analytic version.} Inside the real-analytic Lorentzian category
  $\mathcal C_d^{\mathrm{Lor},\omega}$ there is the subcategory
  $\mathcal C_d^{\mathrm{gh},\omega} \subseteq \mathcal C_d^{\mathrm{Lor},\omega}$ of
  globally hyperbolic cobordisms, in which the manifolds, their Lorentzian metrics,
  and the time-function are all real-analytic. The real-analyticity of the
  time-function can be dropped, since any smooth function can be approximated
  real-analytically. This is the category required for the operator assignment of
  Section~5: the forthcoming Principle~5.1 and Theorem~5.2 move a cobordism inside a
  complexification $M_{\mathbb C}$, a holomorphic thickening that exists only when the
  cobordism and its metric are real-analytic.

  \emph{The metric normal form and the deformation into allowable metrics.} For a
  globally hyperbolic cobordism equipped with a time-function, the metric, written
  through the diffeomorphism $M \cong \Sigma_0 \times [0,1]$, takes the form
  \[ h_t + c^2 \, \mathrm{d}t^2, \qquad c \colon \Sigma_0 \times [0,1] \to \mathrm{i}\mathbb R, \]
  where $t \mapsto h_t$ is a family of Riemannian metrics on the time-slices. The
  factor $c$ is purely imaginary, recording the Lorentzian signature: $c^2$ is real
  and negative, so $h_t + c^2\,\mathrm{d}t^2 = h_t - |c|^2\,\mathrm{d}t^2$. A small
  deformation $\delta c$ pushing $c$ into the open right half-plane
  $\mathbb C_+ = \{z \in \mathbb C : \mathrm{Re}\,z > 0\}$ turns the Lorentzian metric
  into an \emph{allowable} complex metric in the sense of Section~2
  (Definition~\ref{def:allowable}): the single eigenvalue $c^2$, which for the
  Lorentzian metric sat on the negative real axis, moves off it into the complex
  plane, while the $h_t$ eigenvalues stay positive, so the angle condition holds for
  small $\delta c$. This is the deformation picture connecting
  $\mathcal C_d^{\mathrm{gh}}$ to the interior $\mathcal C_d^{\mathbb C}$. One would
  like to define the operator $Z_M$ of a Lorentzian cobordism as a limit of the
  operators attached to such deformations, but the deformed metric depends on the
  choice of time-function, which should be irrelevant to $Z_M$. Resolving this is the
  purpose of the forthcoming Principle~5.1 and Theorem~5.2.

  Stating this in Lean requires Lorentzian-manifold geometry not developed here:
  timelike geodesics and their maximal extension, the light-cone structure of a
  Lorentzian metric and the notion of the positive light-cone, smooth time-functions
  and their gradients, the condition that a gradient lie in the positive light-cone,
  fibrations with Riemannian fibres and the orthogonal trajectories giving
  $M \cong \Sigma_0 \times [0,1]$, the metric normal forms $h_t - \mathrm{d}t^2$ and
  $h_t + c^2\,\mathrm{d}t^2$, the openness of the time-function condition, and, for
  $\mathcal C_d^{\mathrm{gh},\omega}$, the real-analytic structure on the manifolds,
  metrics, and time-function together with the complexification $M_{\mathbb C}$.
\end{definition}

\section{Wick rotation of a time-symmetric germ}

\begin{definition}[Wick rotation of a time-symmetric germ; KS Section 5]
  \label{def:wick-rotation}
  \notready
  \uses{def:lorentzian-cobordism-category, def:unitarity, def:field-theory}
  For a \emph{time-symmetric} Lorentzian germ the Wick rotation continues the
  Lorentzian metric to a Riemannian one and uses it to define the Euclidean space
  $E_{\Sigma}$.

  \emph{The setup.} Let $\Sigma \subset U$ be a Lorentzian germ, with $\Sigma$
  co-oriented in $U$. Exponentiating the geodesics emanating perpendicularly from
  $\Sigma$ identifies $U$ with $\Sigma \times (-\varepsilon, \varepsilon)$, and in
  this identification the Lorentzian metric takes the geodesic-normal form
  \[ h_t - \mathrm{d}t^2, \]
  where $t \mapsto h_t$ is a smooth map from $(-\varepsilon, \varepsilon)$ into the
  manifold of Riemannian metrics on $\Sigma$, as in
  Definition~\ref{def:lorentzian-cobordism-category}.

  \emph{The Wick rotation.} Suppose the germ is \emph{time-symmetric} in the sense
  of Definition~\ref{def:unitarity}, that is $\Sigma \cong \bar\Sigma^*$. We define
  the Euclidean space $E_{\Sigma}$ of the Lorentzian germ by replacing its
  Lorentzian metric with the \emph{Wick-rotated} Riemannian metric
  \[ h_{\mathrm{i}t} + \mathrm{d}t^2, \]
  and taking $E_{\Sigma}$ to be the space that the field theory of
  Definition~\ref{def:field-theory} assigns to this Riemannian germ.

  \emph{Why time-symmetry licenses the rotation.} For a Lorentzian germ in normal
  form the time-symmetry condition $\Sigma \cong \bar\Sigma^*$ is the reflection
  $t \mapsto -t$, which reverses the co-orientation and, the metric being real,
  leaves it unchanged; invariance under the reflection is the relation
  $h_t = h_{-t}$. Thus $t \mapsto h_t$ is even, hence a function of $t^2$: there is
  a smooth $s \mapsto H(s)$ near $0$ with $h_t = H(t^2)$. The rotation
  $t \mapsto \mathrm{i}t$ replaces $t^2$ by $(\mathrm{i}t)^2 = -t^2$, so
  $h_{\mathrm{i}t} = H(-t^2)$ is defined and real, a genuine Riemannian metric on
  $\Sigma$. The rotated metric $h_{\mathrm{i}t} + \mathrm{d}t^2$ is therefore
  Riemannian, of Euclidean signature, and $E_{\Sigma}$ is well defined. Without
  time-symmetry the family $t \mapsto h_t$ need not be even, $h_{\mathrm{i}t}$ need
  not be real, and the construction does not apply.

  \emph{The limitation.} This Wick rotation defines $E_{\Sigma}$ only for
  time-symmetric germs. For a general hypersurface it does not help, and even where
  it applies the prescription is rather arbitrary, tied to the chosen normal form.
  The definition of $E_{\Sigma}$ for a general Lorentzian germ, and the independence
  of the construction from the choices made, are recovered later by the forthcoming
  Principle~5.1 and Remark~5.3.

  Stating this in Lean requires Lorentzian-manifold geometry not developed here:
  the geodesic-normal-form identification
  $U \cong \Sigma \times (-\varepsilon, \varepsilon)$ of a Lorentzian germ, the
  manifold of Riemannian metrics on $\Sigma$ together with the smooth family
  $t \mapsto h_t$, the reflection structure of a time-symmetric germ giving
  $h_t = h_{-t}$, and the analytic continuation $t \mapsto h_{\mathrm{i}t}$, that is
  the passage from $h_t = H(t^2)$ to $H(-t^2)$ which realizes the rotated metric as
  real.
\end{definition}
